\section{投资问题：恒逸的差异化能否穿越周期}

恒逸的竞争优势来自“海外炼化支点+国内完整聚酯链”，而不是任何单一产品产能第一。核心判断是，文莱一期在东南亚供给偏紧时具有区域稀缺性，国内PTA和聚酯规模又提供链条协同；但与荣盛、恒力、东方盛虹相比，文莱上游绝对规模较小，且70\%权益、高杠杆和二期资本开支削弱归母弹性。与桐昆、新凤鸣相比，恒逸拥有上游利润池，却承受更复杂的跨境和资本结构风险。

全球油品环境呈现需求偏弱与产品紧张并存。IEA公开摘要指向炼厂加工减少和裂解利润高位；EIA则预计Brent在2026H2回落。供给收缩能够抬升炼厂利润，却不等于终端需求长期繁荣。这种环境最有利于能稳定运行、产品结构匹配区域缺口且原料物流可控的炼厂，也意味着利润对供给恢复高度敏感。

\section{六家公司不是一组静态PE可比}

荣盛、恒力和东方盛虹的上游规模显著大于恒逸；桐昆和新凤鸣更接近国内PTA--长丝加工模式。不同公司在并表／权益法、少数股东、油煤气路线、港口公用工程和新材料业务上的差异，使静态PE比较容易失真。可比公司的正确用途是校验6--9倍周期PE是否合理，并判断恒逸的区位与资本折价，而不是机械选出最低PE。

\begin{exhibitbox}[一体化竞争格局]
\scriptsize
\noindent\begin{tabularx}{\textwidth}{L{1.75cm}L{2.35cm}L{2.35cm}L{2.35cm}X}
\toprule
\textbf{公司} & \textbf{炼化／上游} & \textbf{PTA／聚酯} & \textbf{相对优势} & \textbf{可比风险} \\
\midrule
恒逸石化 & 文莱一期800万吨；70\%权益 & PTA参控股2150；PET Fiber 938万吨 & 海外区域支点、链条完整 & 上游规模较小、少数股东、杠杆与二期 \\
荣盛石化 & 浙石化4000万吨；PX 1040万吨 & PTA 2150、瓶片530万吨等 & 上游规模、港口管道与材料延伸 & 少数股东与业务体量显著不同 \\
恒力石化 & 原油2000、原煤600万吨 & PTA 1660、民用长丝410万吨 & 油煤化、公用工程与码头自配 & 设计利用率不能外推恒逸 \\
东方盛虹 & 炼化1600；MTO／PDH多路线 & 长丝约360、再生约60万吨 & 原料路线与新能源材料多元 & 新材料使盈利可比性较弱 \\
桐昆股份 & 原油加工权益量1000万吨 & PTA 1020、长丝1510万吨 & 长丝规模和市场份额领先 & 炼化为权益量，非并表炼厂 \\
新凤鸣 & 无大型炼化 & PTA 1100、长丝885、短纤120万吨 & 纯制造链、装置效率可比 & 缺少油品高景气利润池 \\
\bottomrule
\end{tabularx}
\end{exhibitbox}
\sourcenote{HY-OFF-2025AR；IND-PEERS对应五家2025年报；analysis/competitive\_landscape.md，数据截止2025-12-31，研究截止2026-07-22。设计产能、权益量与并表口径不同。}

表中的“优势”只有在盈利和现金流中体现才有估值意义。荣盛与恒力的大规模公用工程、码头和上游自给形成成本锚；桐昆与新凤鸣提供国内聚酯加工效率锚；恒逸的文莱区位在2026年产品紧张时更有弹性，却也更依赖海外单一项目和税惠。

\section{恒逸的四项优势与四项折损}

第一，区域稀缺。文莱一期约占公司引用的东南亚炼油产能3\%，产品直接面对净进口市场。第二，原料采购半径相对短，文莱和马来西亚来源有助于降低相对中国长航线的物流暴露。第三，链条完整，PX、PTA、聚酯和CPL--PA6使利润池有迁移空间。第四，装置与管理经验已经通过一期运营得到验证。

对应的折损也很明确。第一，800万吨／年上游规模小于国内三家大炼化，规模摊薄并非绝对领先。第二，30\%文莱利润自动归少数股东。第三，2025年国内PTA和涤纶毛利率仍薄，无法无条件对冲文莱回落。第四，资产负债率约72\%、担保和二期投入限制资本回报。

\begin{exhibitbox}[竞争力评分卡：从能力到每股价值]
\noindent\begin{tabularx}{\textwidth}{L{2.6cm}C{1.5cm}L{3.45cm}X}
\toprule
\textbf{维度} & \textbf{判断} & \textbf{正面证据} & \textbf{估值折损} \\
\midrule
区域与产品位置 & 强 & 东南亚成品油缺口、油品与芳烃组合 & 价差高位可逆、客户未具名 \\
一体化与采购 & 中强 & 文莱到国内PTA／聚酯双链 & 内部交易和库存占用复杂 \\
规模与成本 & 中 & 多板块产能领先、差异化纤维占比提升 & 上游规模小于国内大炼化 \\
利润归属 & 中弱 & 70\%控股提供控制 & 30\%少数股东分流 \\
资产负债表 & 弱 & 授信与项目经验支持运营 & 高负债、担保、资本开支与转债 \\
证据透明度 & 中弱 & 年报产能和分部较完整 & 2026分部利润、价差、税率未披露 \\
\bottomrule
\end{tabularx}
\end{exhibitbox}
\sourcenote{HY-OFF-2025AR；HY-OFF-2026Q1；analysis/competitive\_landscape.md；analysis/company\_fundamental\_cards.md，研究截止2026-07-22。}

\section{集中度提升不保证加工费}

聚酯长丝集中度已经较高，桐昆年报引用协会数据称六家上市公司直纺长丝产能占比约81\%。集中度有助于检修协调和行业自律，却没有阻止2025年PTA加工费处于近五年低位，也没有阻止聚酯利润下降。商品行业的集中度只有在供给纪律、库存和终端需求共同改善时才形成利润。

因此，恒逸的国内竞争优势不应写成“强者恒强”，而应落到装置效率、差异化产品、库存周转和现金成本。公司称差别化纤维产量占比提升至35\%，这支持产品结构升级方向；但缺少差异化产品收入、溢价和毛利拆分，不能直接转为估值溢价。

\begin{keyinsight}[竞争格局的“所以呢”]
文莱区位使恒逸在2026年成为利润弹性标的，但高弹性不等于高确定性。\textbf{行动上，估值只给予7.5倍基准周期PE；要向9倍牛情景升级，必须看到区域价差持续、国内库存下降、现金流覆盖资本开支和二期融资受控。}
\end{keyinsight}

\section{竞争地位的证伪与升级}

竞争优势失效的证据包括：东南亚炼厂复产使产品裂解快速收窄、文莱检修或运输中断、PTA复产后再累库、聚酯库存和出口同时恶化，以及二期扩张削弱一期现金回流。上述任一因素都可能让“海外区位”从溢价来源变成风险集中。

升级证据则包括：文莱季度分部利润、税率与产品价差可核验；国内PTA与聚酯分装置开工和单吨加工费改善；差异化产品溢价被财务数据确认；二期贷款和资本金在明确回报门槛下交割。在此之前，同行只作卫星验证，不发布额外目标价，也不把横向产能排名转成推荐排序。

\paragraph{最强反方情景} 若国内大炼化凭借更大规模、公用工程和物流继续压低成本，而东南亚区域溢价消失，恒逸将同时面对上游规模劣势和国内低毛利资产回报。\textbf{研究上，这一组合出现时，6倍PE熊情景应取代7.5倍基准。}
